% ======================================================================
% Kapitel 5 — Conclusion, Limitations and Outlook (Budget: ~2 Seiten)
% Schlank halten: KEINE erneute Ergebnisdiskussion (die steht in Kap. 4).
% ======================================================================
\chapter{Conclusion, limitations and outlook}
\label{ch:conclusion}

\section{Summary and contributions}
\label{sec:summary}
% TODO: Was wurde gebaut und gezeigt (3--4 Absaetze); wissenschaftlicher
% und praktischer Beitrag explizit benennen (vgl. Salma 5.1).
...

\section{Limitations}
\label{sec:limitations}
% TODO: methodische Einschraenkungen als Bullets (vgl. Salma 5.2):
% Korpusgroesse/-auswahl, Query-Set-Umfang (28), LLM-as-judge-Bias
% (+ Beleg: Non-Committal-Klassifikator traf 8B UND Haiku ---
% modell-agnostische Judge-Metrik-Kritik),
% Query-Set nennt prominente CVEs -> beguenstigt parametrische Baseline
% systematisch,
% Selbstevaluierungs-Bias (szenariobasierte Bewertung durch den Autor,
% mitigiert durch das vorab definierte Raster in Anhang C),
% keine Nutzerstudie, Einzelsprache, C3 ohne Laufzeit-Credibility-Scoring,
% Context Precision/Recall sind referenzbasierte, LLM-gejudgte Proxys,
% KEINE klassischen IR-Metriken; Qrels-basierte Ranking-Evaluation bleibt
% Future Work (per-Stage-Retrieval-Traces liegen in den Artefakten vor), ...
...

\section{Future work}
\label{sec:future-work}
% TODO: zweigeteilt nach Salmas Vorbild:
% (a) praktische Roadmap (z.B. MCP-/Agent-Integration, Feed-Anbindung,
%     Multi-User-Betrieb),
% (b) wissenschaftliche Anschlussmoeglichkeiten (Nutzerstudie mit Analysten,
%     groessere Korpora, weitere Modellvergleiche, Qrels-basierte
%     IR-Evaluation auf Basis der vorhandenen Retrieval-Traces).
...

% TODO Abbildung: Roadmap-Kaskade nach Salmas Vorbild (Abb. 8 bei ihr):
% gestufte Pfeile MCP-Integration -> Feed-Anbindung -> Multi-User -> Nutzerstudie.
% \begin{figure}[htbp]
%     \centering
%     \includegraphics[width=0.85\textwidth]{roadmap_outlook}
%     \caption{Practical roadmap towards operational deployment}
%     \label{fig:roadmap}
% \end{figure}
